\documentclass[10pt,a4paper]{amsart}
\usepackage{amsmath,amssymb,amsthm,graphicx}
% For commutative diagrams we use Xy-pic
\usepackage[all]{xy}

% We use 2cell for 2-commutative diagrams.
\xyoption{2cell}
% Theorem environments.
%
\theoremstyle{plain}
\newtheorem{theorem}[subsection]{Theorem}
\newtheorem{proposition}[subsection]{Proposition}
\newtheorem{lemma}[subsection]{Lemma}

\theoremstyle{definition}
\newtheorem{definition}[subsection]{Definition}
\newtheorem{example}[subsection]{Example}
\newtheorem{exercise}[subsection]{Exercise}
\newtheorem{situation}[subsection]{Situation}

\theoremstyle{remark}
\newtheorem{remark}[subsection]{Remark}
\newtheorem{remarks}[subsection]{Remarks}

\numberwithin{equation}{subsection}

% Macros
%
\def\lim{\mathop{\mathrm{lim}}\nolimits}
\def\colim{\mathop{\mathrm{colim}}\nolimits}
\def\Spec{\mathop{\mathrm{Spec}}}
\def\Hom{\mathop{\mathrm{Hom}}\nolimits}
\def\Ext{\mathop{\mathrm{Ext}}\nolimits}
\def\SheafHom{\mathop{\mathcal{H}\!\mathit{om}}\nolimits}
\def\SheafExt{\mathop{\mathcal{E}\!\mathit{xt}}\nolimits}
\def\Sch{\mathit{Sch}}
\def\Mor{\mathop{\mathrm{Mor}}\nolimits}
\def\Ob{\mathop{\mathrm{Ob}}\nolimits}
\def\Sh{\mathop{\mathit{Sh}}\nolimits}
\def\NL{\mathop{N\!L}\nolimits}
\def\CH{\mathop{\mathrm{CH}}\nolimits}
\def\proetale{{pro\text{-}\acute{e}tale}}
\def\etale{{\acute{e}tale}}
\def\QCoh{\mathit{QCoh}}
\def\Ker{\mathop{\mathrm{Ker}}}
\def\Im{\mathop{\mathrm{Im}}}
\def\Coker{\mathop{\mathrm{Coker}}}
\def\Coim{\mathop{\mathrm{Coim}}}

% Boxtimes
%
\DeclareMathSymbol{\boxtimes}{\mathbin}{AMSa}{"02}

%
% Macros for moduli stacks/spaces
%
\def\QCohstack{\mathcal{QC}\!\mathit{oh}}
\def\Cohstack{\mathcal{C}\!\mathit{oh}}
\def\Spacesstack{\mathcal{S}\!\mathit{paces}}
\def\Quotfunctor{\mathrm{Quot}}
\def\Hilbfunctor{\mathrm{Hilb}}
\def\Curvesstack{\mathcal{C}\!\mathit{urves}}
\def\Polarizedstack{\mathcal{P}\!\mathit{olarized}}
\def\Complexesstack{\mathcal{C}\!\mathit{omplexes}}
% \Pic is the operator that assigns to X its picard group, usage \Pic(X)
% \Picardstack_{X/B} denotes the Picard stack of X over B
% \Picardfunctor_{X/B} denotes the Picard functor of X over B
\def\Pic{\mathop{\mathrm{Pic}}\nolimits}
\def\Picardstack{\mathcal{P}\!\mathit{ic}}
\def\Picardfunctor{\mathrm{Pic}}
\def\Deformationcategory{\mathcal{D}\!\mathit{ef}}
\usepackage{cleveref}
\crefname{originallemma}{original lemma}{original lemmas}
\Crefname{originallemma}{Original lemma}{Original lemmas}
\makeatletter
\newlabel{lemma-lift-formal-smoothness}{{031L}{source}{}{sourcepointer-1}{}}
\newlabel{lemma-lift-formal-smoothness@cref}{{[originallemma][0][]031L}{[1][0][]source}{}{}{}}
\makeatother
\crefname{originallemma}{original lemma}{original lemmas}
\Crefname{originallemma}{Original lemma}{Original lemmas}
\makeatletter
\newlabel{proposition-characterize-separable-field-extensions}{{0322}{source}{}{sourcepointer-2}{}}
\newlabel{proposition-characterize-separable-field-extensions@cref}{{[originallemma][0][]0322}{[1][0][]source}{}{}{}}
\makeatother
\title{The reduction of a Cohen ring modulo a power of the residue characteristic is formally smooth over Z/p\textasciicircum{}n}
\author{The Stacks Project Authors}
\date{stacks-project:algebra.tex:local-snapshot-20261012; GNU FDL 1.2 or later}
\begin{document}
\maketitle
\noindent\textit{Editorial note.} All delivered mathematical text is the author's own, copied from the pinned Stacks Project chapter file algebra.tex; no statement or proof step was rewritten, reordered or removed.\par\medskip
\noindent\textit{Editorial note.} Preamble provenance, as required for this source: the first line of the pinned chapter file is \textbackslash{}input\{preamble\}, and the native preamble was not staged with the source, so no preamble text lives inside the chapter. The retained preamble used here is the project's single Stacks-wide preamble (the same 217-line / 8116-byte file that the already published chow and etale-cohomology cards use), copied mechanically line for line from sources/stacks-chow/preamble.tex lines 14-18 (Xy-pic setup) and 154-217 (theorem declarations and mathematical macros). Not one character was changed, no package was added, no macro was defined, and the excerpt serves only as preamble, never as a span coordinate.\par\medskip
\noindent\textit{Editorial note.} The same-chapter pointers are located at their \textbackslash{}label lines inside the pinned chapter file; the cross-chapter pointers are located in the project's own tag list staged in the snapshot (sources/stacks-chow/tags.txt), at the verbatim lines carrying their canonical tags.\par\medskip
\noindent\textit{Editorial note.} The staged candidate list records ext\_refs = 0 for this label, but the proof actually references two same-chapter results, which the generator resolved and recorded as navigation stubs; the induction itself stays inside the counted span. This is the second instance today of that staged field being unreliable.\par\medskip
\noindent\textit{Editorial note.} The counted proof contains no bibliography citation, so no bibliography excerpt is attached.\par\medskip
\noindent\textit{Allowed standard prerequisites (editorial).}\begin{itemize}
\item Characterisation of separable field extensions by formal smoothness. Same chapter, Proposition labelled proposition-characterize-separable-field-extensions; invoked at proof line 45135. Settles the case n = 1, since the residue field of a Cohen ring of characteristic p is separable over its prime field.
\item Lifting formal smoothness along a square-zero ideal. Same chapter, Lemma labelled lemma-lift-formal-smoothness; invoked at proof lines 45138-45141. Performs the induction step, from level n to level n + 1, for the ideal generated by p\textasciicircum{}n.
\end{itemize}\medskip
\ifdefined\hypertarget\hypertarget{sourcepointer-1}{}\fi
\noindent\textit{Source reference.} Original lemma 031L: sources/stacks-algebra/pinned-original-source.tex, source lines 38317--38317. Navigation pointer for \textbackslash{}ref\{lemma-lift-formal-smoothness\} used inside the counted span (line 45139); the argument that consumes it is retained in the span.\par\medskip
\ifdefined\hypertarget\hypertarget{sourcepointer-2}{}\fi
\noindent\textit{Source reference.} Original lemma 0322: sources/stacks-algebra/pinned-original-source.tex, source lines 44692--44692. Navigation pointer for \textbackslash{}ref\{proposition-characterize-separable-field-extensions\} used inside the counted span (line 45136); the argument that consumes it is retained in the span.\par\medskip
\begin{lemma}
\label{lemma-cohen-ring-formally-smooth}
Let $p > 0$ be a prime.
Let $\Lambda$ be a Cohen ring with residue field of characteristic $p$.
For every $n \geq 1$ the ring map
$$
\mathbf{Z}/p^n\mathbf{Z} \to \Lambda/p^n\Lambda
$$
is formally smooth.
\end{lemma}

\begin{proof}
If $n = 1$, this follows from
Proposition \ref{proposition-characterize-separable-field-extensions}.
For general $n$ we argue by induction on $n$.
Namely, if $\mathbf{Z}/p^n\mathbf{Z} \to \Lambda/p^n\Lambda$ is
formally smooth, then we can apply Lemma \ref{lemma-lift-formal-smoothness}
to the ring map
$\mathbf{Z}/p^{n + 1}\mathbf{Z} \to \Lambda/p^{n + 1}\Lambda$
and the ideal $I = (p^n) \subset \mathbf{Z}/p^{n + 1}\mathbf{Z}$.
\end{proof}

\end{document}
